\section{格点研究的数值细节}

前一节中我们描述了处理算符连续极限的存在性问题与大 $N$ 相变点问题的策略。我们也努力使自己相信大 $N$ 相变确实是连续的。现在我们着手描述模拟的准确细节以及得到的数字。

\subsection{$f_c(l)$ 的确定}

\subsubsection{估计无穷大 $N$ 的能隙}

在无穷大 $N$，小的 $\hat W$ 圈在格点上会有一个谱隙。我们需要一个方法从有限 $N$ 的数值数据中提取该能隙。当我们胀大圈和/或减少涂抹量时，能隙收缩。我们需要一种外推到能隙消失点的方法。我们希望避开对太小的能隙作测量，只通过外推进入那一区域。

获得能隙的无穷大 $N$ 值是困难的，因为在能隙边缘次领头效应很大。合理的是假定能隙边缘的普适行为与简单厄米矩阵模型中能隙边缘的行为相同。在这个假设下，对本征值密度的行为和极端本征值的行为我们有分开的预言。可以数值地检查普适形式工作得多好，我们发现它们工作得相当好。为了拟合到普适行为上，我们需要调节一个对应于无穷大 $N$ 能隙大小的参数。这就给出了我们随后使用的能隙大小的测定。

只有在本征值密度有能隙的相变那一侧，才存在大 $N$ 修正被普适函数主导的区域；因此我们能够用单个（但很大的）$N$ 得到的数据获得无穷大 $N$ 本征值能隙大小的良好估计。为稳妥起见，我们检验了所预言的对 $N$ 的有理幂次依赖在我们分析的区域内成立。与一些 $N=28$ 模拟的比较使我们相信，在 $N=28,44$ 处对 $N$ 的幂次依赖已达到已知的
普适值。

我们在固定的较大 $N$ 下研究了 $\hat W [L , L ; f;
n=L^2;N]$ 在格点上对若干组 $L$ 和 $b$ 值的本征值分布。
令 $e^{i\theta_k}, k=1,\cdots,N$ 为格点上固定 $L\times L$ 圈的 $N$ 个本征值，$-\pi \le \theta_1 <
\theta_2 < \cdots < \theta_{N-1} < \theta_N < \pi$。
所有角度之和模 $2\pi$ 为零。在有 $V$
个格点的格点上存在 $6V$ 个固定取向的不同 $L\times L$ 威尔逊圈；我们计算了所有这类圈的本征值。
以此方式我们得到固定规范场背景 $U$ 中所有本征值的分布
$p_f(\theta;U)$ 的直方图。
我们还收集了固定规范场背景中最大本征值分布
$p_m(\theta_N;U)$ 的直方图数据。此外，我们记录了关于最大本征值的均值、方差、偏度和峰度：
\begin{eqnarray}
\mu_N(U) &=& \langle \theta_N \rangle_U;\nonumber\\
\sigma^2_N(U) &=& \langle \theta^2_N \rangle_U - \langle
\theta_N \rangle^2_U;\nonumber\\[0.6em]
S_N(U) &=&
\frac{ \langle \theta^3_N \rangle_U
-3 \langle \theta^2_N \rangle_U \langle \theta_N \rangle_U
+ 2 \langle \theta_N \rangle^3_U }{\sigma^3_N(U)};\nonumber\\[0.6em]
K_N(U) &=&
\frac{ \langle \theta^4_N \rangle_U
-4 \langle \theta^3_N \rangle_U \langle \theta_N \rangle_U
+6 \langle \theta^2_N \rangle_U \langle \theta_N \rangle^2_U
-3 \langle \theta_N \rangle^4_U }{\sigma^4_N(U)}-3.
\end{eqnarray}

如果 $\hat W$ 的本征值分布有可观的能隙，那么忽略两条边缘相互之间的影响是合理的。于是能隙边缘的普适行为应当与简单高斯厄米矩阵模型中能隙边缘的行为相同。在此假设下，我们把数据与极端本征值的行为~\cite{tracy}
以及边缘附近本征值密度的行为~\cite{gff}作了比较。

极端本征值的普适分布 $p(s)$ 由文献~\cite{tracy} 给出：
\be
p(s)=\left [ \int_s^\infty q^2(x) dx \right ] e^{-\int_s^\infty (x-s)q^2(x) dx}
\ee
其中 $q$ 是 Painleve II 方程的解，
\be
\frac{d^2q}{ds^2} = sq + 2 q^3;
\ee
并满足渐近条件
\be
q(s) \sim Ai(s)\ \ {\rm as}\ \ s\rightarrow\infty.
\ee
上述分布的均值、方差、偏度和峰度分别为 $-1.7710868074$、$0.8131947928$、
$0.2240842036$ 和 $0.0934480876$。为了把 $\theta_N$ 分布的格点数据匹配到上述普适形式，
我们按下式定义 $s$：
\be
\theta_N = \frac{s}{\alpha} + \pi(1-g)
\label{scalevar}
\ee
并令
\begin{eqnarray}
\alpha &=& \sqrt{\frac{0.8131947928}{\sigma^2_N}}\nonumber\\
g &=& 1-\frac{1}{\pi}[\mu_N + \frac{1.7710868074}{\alpha}]
\end{eqnarray}
$\alpha$ 是非普适的标度因子，$g$ 是单位圆上以 $-1$ 为中心的本征值分布中能隙弧长的一半。量 $\sigma^2_N$ 与 $\mu_N$
是前面定义的量 $\sigma^2_N(U)$ 与 $\mu_N(U)$ 在规范场系综上的平均。从 $\sigma^2_N(U)$ 与 $\mu_N(U)$ 的格点数据，我们得到 $\alpha$ 与 $g$ 的估计及用 jackknife（刀切）方法提取的误差。利用这些估计，我们对固定的 $f$, $b$, $N$, $V$ 把标度化后的格点数据与 $p(s)$ 比较。其中一个比较示于图~\ref{gap}：$f=0.15$, $b=0.361653$, $L=4$, $N=44$, $V=7^4$，能隙为 $0.1229(4)$，两者符合良好。尽管需要整个 $s$ 范围的数据才能得到偏度和峰度的良好估计，我们发现对很宽的能隙值范围，格点数据与普适估计在偏度上符合到 $5\%$ 以内、峰度上符合到 $20\%$ 以内。偏度与峰度的估计不依赖于提取出的 $\alpha$ 和 $g$ 值。

\begin{figure}[ht]
\centering
\includegraphics[width=\textwidth]{images/gap.pdf}
\caption{格点上得到的标度化极端本征值分布（圆点）（取 $f=0.15$, $b=0.361653$, $L=4$, $N=44$, $V=7^4$）与普适分布（实曲线）的比较。}
\label{gap}
\end{figure}

上述数值分析的主要结果是固定的 $f, L, b$ 下的半能隙大小 $g(f,L,b)$。


还有另一种获得 $g$ 估计的方法：
最近的一个结果给出了高斯厄米矩阵本征值密度的渐近修正~\cite{gff}。利用该公式，我们用 (\ref{scalevar}) 中的标度化变量把本征值密度拟合到
\begin{equation}
p_f(s) = c \left\{ [Ai^\prime(s)]^2 - s [Ai(s)]^2 \right\} +d
\left\{ 3s^2[Ai(s)]^2 - 2s [Ai^\prime(s)]^2 -3Ai(s)Ai^\prime(s)
\right\} \label{airy}
\end{equation}
变量 $c$ 由整体归一化固定；在渐近公式~\cite{gff} 中 $\frac{d}{c}=-\frac{1}{20N^{2/3}}$，其中 $N$ 是矩阵的大小。
参数 $d$ 可以是非普适的；$c$ 在某种意义上也可以是非普适的，因为它依赖于本征值的总数，而且高斯厄米矩阵模型是否以相同的 $N$ 匹配到我们的模型并不清楚。因此，除 (\ref{scalevar}) 中出现的两个变量之外，我们还把 $c$ 和 $d$ 作为独立变量把格点分布拟合到上述渐近公式，产生了 $g(f,L,b)$ 的第二次测定。

\begin{figure}[ht]
\centering
\includegraphics[width=\textwidth]{images/dist.pdf}
\caption{格点上得到的本征值分布（带虚线的圆点）（取 $f=0.15$,
$b=0.361653$, $L=4$, $N=44$, $V=7^4$）与边缘附近普适分布（实曲线）的比较。 }
\label{dist}
\end{figure}

图~\ref{dist} 显示了完整本征值分布与边缘附近普适分布的比较。普适分布抓住了尾部，其原因与极端本征值分布符合 $p(s)$ 相同。但我们也看到与边缘附近头两个振荡的良好符合，这包含了超出极端本征值的内容。总的来说，我们发现能隙的测定与只用极端本征值所得的结果一致。误差估计我们只用了极端本征值方法。



\subsubsection{外推到零能隙}



下一步是提取固定的 $b$ 和 $L$ 下无穷大 $N$ 的 $f$ 临界值。做法是假设如下形式的拟合
\be
g(f,L,b) = A(L,b) \sqrt{f - F_c(L,b)}
\ee
对以解析方式进入概率律的外部参数的平方根依赖在矩阵模型中是普遍的，并且得到数据的充分支持。

$b=0.361653$、$N=44$ 处 $g^2(f,L,b)$ 作为 $f$ 函数的图示于图~\ref{gap_vs_ff}。提取 $F_c(L,b)$ 时使用了一段 $f$ 范围，使能隙 $g(f,L,b)$ 既不太接近 $0$ 也不过远离 $0$。如果走到太小的能隙，我们担心未被计及的圆化效应；若能隙太大，就没有什么理由使用我们假定的普适平方根行为。

图~\ref{gap_vs_ff} 还显示了拟合并连同 $F_c(L,b)$ 的估计值。


\begin{figure}[ht]
\centering
\includegraphics[width=\textwidth]{images/gap_vs_ff.pdf}
\caption{$b=0.361653$, $N=44$, $V=7^4$ 处 $g^2(f,L,b)$ 作为 $f$ 函数的图。还显示了对数据的拟合以及外推得到的 $F_c(L,b)$ 值。 }
\label{gap_vs_ff}
\end{figure}


\subsubsection{$F_c(L,b)$ 的连续统极限}


我们现在着手取 $F_c(L,b)$ 的连续统极限。
当改变 $b$ 时，我们用两圈蝌蚪改进微扰论保持物理圈固定。
利用~\cite{knn}，我们有以蝌蚪改进微扰论表示的、以格点单位计的临界退禁闭温度的数值公式，其中 $e(b)$ 是耦合 $b$ 处未扭转单方块规范作用量的方块平均值。
\begin{equation}
\frac{1}{T_c(b)} = 0.26
\left[\frac{11}{48\pi^2 b_I}\right]^{\frac{51}{121}}
e^{\frac{24\pi^2 b_I}{11}};\ \ \ \
b_I=be(b)
\label{twoloop}
\end{equation}
这将用于确定物理单位下的格距。

我们设定 $lt_c=LT_c(b)$ 并在若干物理圈尺寸值处积累数据：$lt_c=0.42,0.48,0.58,0.64$。在每个这样的物理尺寸处，我们经历了三到四个格距 $a(b)$ 的序列，并且对每个 $b$ 用上面解释的方法得到一个 $F_c(L=\frac{l}{a(b)},b)$ 值。

对每个 $lt_c$，然后把 $F_c(L=\frac{l}{a(b)},b)$ 值外推到连续统，假设
\be
F_c(L=\frac{l}{a(b)},b) = f_c(lt_c) + O(T_c^2(b)).
\ee
这里回顾 $T_c(b)=a(b)t_c$。（我们将对同一量使用略欠一致的记号 $f_c(l)$ 和 $f_c(lt_c)$：其宗量或有量纲，或用物理退禁闭温度 $t_c$ 的适当幂次化为无量纲。）

$f_c(lt_c)$ 的外推结果示于图~\ref{crit2}。


\begin{figure}[ht]
\centering
\includegraphics[width=\textwidth]{images/crit2.pdf}
\caption{$F_c(L=\frac{l}{a(b)},b)$ 作为 $T_c(b)$ 函数的图，对应四个不同的 $lt_c$ 值。还显示了数据的拟合，给出 $f_c(lt_c)$ 的连续统值。 }
\label{crit2}
\end{figure}



我们的最后结果绘在 $(f,l)$ 平面中。临界线 $f_c(l)\equiv f_c (lt_c)$ 把该平面分成两部分，如图~\ref{fltc} 所示。上半部分中威尔逊圈的本征值分布有能隙，下半部分则无隙。以零格距临界线为下界的带状区域指示有限格距效应的范围。



\subsection{标度变换下相变两侧的研究}

图~\ref{fltc} 中有限格点尺寸效应明显。仅仅去标度圈并不现实，因为当 $l$ 在合适的范围内变化时曲线 $f_c(l)$ 的变化小于格点修正。要看清标度的效应，我们需要走 $(f,l)$ 平面中一条非水平的线，这意味着我们将把涂抹作为物理圈尺寸的函数来改变。照例，我们把 $l$ 换成无量纲变量 $lt_c$。

我们选了这条线（下文称之为``穿越线''）
\begin{equation}
f = 0.25 - (0.6128 lt_c)^2
\label{fline}
\end{equation}
如图~\ref{fltc} 所示，它横穿由临界线 $f_c(lt_c)$ 所界的整个有限格点修正带。


考虑格点上尺寸为 $L$ 的固定圈，对各种格点耦合 $b$，使 $f(b)$ 由 (\ref{fline}) 给出，其中 $lt_c=LT_c(b)$ 且 $T_c(b)$ 由 (\ref{twoloop}) 给出。我们希望沿线研究本征值分布中相变的性质。

\begin{figure}[ht]
\centering
\includegraphics[width=\textwidth]{images/fltc.pdf}
\caption{$(f,lt_c)$ 平面中的临界线。还显示了用于研究胀大威尔逊圈所诱导相变的穿越线。 }
\label{fltc}
\end{figure}

我们已经知道，分布对 $l > l_c$ 无能隙而对 $l < l_c$ 有能隙。除了得到 $l_ct_c$ 的估计之外，我们还可以通过固定 $L$ 并沿 (\ref{fline}) 给出的穿越线行进来研究这个相变的普适类。我们取 $L=3,4,5$，并在下列 $b$ 范围内工作：L=3 时 $0.345 < b < 0.365$，L=4 时 $0.346 < b <
0.368$，L=5 时 $0.351 < b < 0.369$。在每个范围内，随着 $b$ 减小物理圈被放大。这样就在同一物理圈尺寸处得到交叠：一次由较小的 $L$ 和较小的 $b$ 代表，另一次由相应较大的这些参数代表。由于标度破坏，不会得到完全相同的本征值分布；我们需要看到，随着 $b$ 和 $L$ 增大这些修正减小，并且逼近一个有限的极限本征值分布。这就是与尺寸 $l$ 的连续统圈相联系的本征值分布。

在图~\ref{fltc} 中，我们对每个 $L=3,4,5$ 显示了从每一相独立外推得到的两个临界长度，它们表现为穿越线上的三角形。随着 $L$ 增大，我们看到每对中两个成员之间的距离随 $L$ 增大而减小。而且，连接每对两点的线段的中点似乎随 $L$ 增大而收敛。这表明在连续统极限中，这条线上将只有一个单一的相变点，即 ${\hat
W}$ 的大 $N$ 谱发生连续相变的一个圈尺寸。

\subsubsection{从小圈一侧估计临界尺寸}\label{gapside}

我们再次用极端本征值分布（$p(s)$）和边缘附近的本征值密度（$p_f(s)$）得到无穷大 $N$ 能隙的估计。两者给出的能隙估计彼此一致，结果绘于图~\ref{gap_diag}。圈的临界尺寸仍通过假设写成如下形式的平方根奇点来提取：
\be
g^2(lt_c)= A \ln \frac{l}{l_c} \ee
对每个 $L=3,4,5$ 我们得到一个 $l_ct_c$ 的估计，显示在图~\ref{gap_diag} 中。可以看到线 $L=4,5$ 彼此比线 $L=3,4$ 更接近；两对线之间近似距离之比大致符合 $a^2 (b)$ 量级修正的假设。

\begin{figure}[ht]
\centering
\includegraphics[width=\textwidth]{images/gap_diag.pdf}
\caption{沿穿越线，能隙的平方作为 $\ln(lt_c)$ 的函数。 }
\label{gap_diag}
\end{figure}


\subsubsection{从大圈一侧估计临界尺寸}\label{rhopi}



对无能隙的算符我们需要一个方法提取角度 $\pi$ 处本征值密度 $\hat\rho
(\pi)$ 的无穷大 $N$ 值。收缩圈和/或加重涂抹使 $\hat\rho
(\pi)$ 减小。我们还需要一种外推到 $\hat\rho(\pi)$ 消失点的方法，而又不过分接近它。

估计 $\hat \rho(\pi)$ 的无穷大 $N$ 值也受到大的次领头效应之害，典型为 $\frac{1}{N}$ 量级。本征值密度中存在反映本征值之间排斥的快速涨落。把数值得到的分箱本征值分布对角度作傅里叶变换、保留相对 $\frac{1}{N}$ 较大的波长，即可消除这种快速涨落。这等价于对相继的 $k$ 值（都保持远小于 $N$）计算 $\hat W^k$ 的迹。只保留那些傅里叶模式，我们就得到 $\hat \rho(\pi)$ 的估计；与较低 $N$ 的模拟所作的核对表明，我们以合理的精度确定 $\hat \rho(\pi)$。


当圈尺寸减小和/或涂抹量增加时把 $\hat \rho(\pi)$ 外推到零的做法与我们处理能隙时类似。我们发现领头行为又是平方根型的，但这只在小区间内成立，可以通过增加更多 $(l-l_c)$ 的幂次来扩展。用这种方法确定的 $l_c$ 相当稳健。

在二维中，任意尺寸圈的大 $N$ 本征值分布在连续统中已知 \cite{duol}。
正如我们稍后将详细阐述的，可以合理地假设靠近大 $N$ 相变的行为是普适的，并且对我们的四维 $\hat \rho(\theta)$ 也成立。
因此，我们假定
${\hat \rho(\pi)}$
在任何以解析方式进入控制涂抹威尔逊圈算符的底层分布的参数的临界值处按平方根消失。这就是为什么当从上方接近相应的临界尺寸 $l_c$ 时，我们取 $\hat\rho (\pi)$ 按
$\sqrt{l-l_c}$ 消失。


如上所述，对相对 $l_c$ 足够大的 $l$，我们先对所有角度把分布表示为：
\be
\hat\rho(\theta) = \sum_{k=0}^N f_k \cos(k\theta)
\ee
然后截断该傅里叶级数，并用截断级数给出 $\hat\rho(\pi)$ 的值：
\be
\hat\rho(\pi) = \sum_{k=0}^{k_m} (-1)^k f_k
\ee
$k_m$ 选得略小于 $N/2$；这消除了频率为 $N$ 的振荡。我们通过观察 $k_m$ 变动时的稳定性核验了这一选择的合理性。

我们
发现，对足够低的 $k$，$f_k$ 从
$k=1$ 起符号振荡。
如图~\ref{fourier} 所示，直到约 $\frac{N}{2}$ 的 $k$，系数对 $k$ 光滑，且符号振荡得到遵守。稍后我们还要回到这个性质。

\begin{figure}[ht]
\centering
\includegraphics[width=\textwidth]{images/fourier.pdf}
\caption{无隙区域中两个不同圈的傅里叶系数 $f_k$ 作为 $k$ 的函数。测量是用 $N=44$ 的 $L=5$ 圈在两个不同格距下获得的。 }
\label{fourier}
\end{figure}



我们把得到的 $\hat\rho(\pi)$ 值在临界区域拟合为平方根，结果显示在图~\ref{rho_diag} 中。
由 $\hat\rho(\pi)$ 得到的 $l_ct_c$ 估计往往高于由拟合能隙得到的估计。这一差别随格距趋于零而减小。

由能隙和由 $\hat\rho(\pi)$ 得到的 $l_ct_c$ 临界值都显示在图~\ref{fltc} 中。可以看到两个估计之间的差别在收缩；还可以看到两个估计的平均值随格距减小而逼近一个有限极限。


我们所看到的标度破坏的大小和行为与通常预期一致，反映了更高维连续统算符对作用量和算符的贡献。对这些标度破坏更详细的研究留待将来，但指出其数量级是典型的。



\begin{figure}[ht]
\centering
\includegraphics[width=\textwidth]{images/rho_diag.pdf}
\caption{沿穿越线，$\hat\rho(\pi)$ 的平方作为 $\ln(lt_c)$ 的函数。 }
\label{rho_diag}
\end{figure}


\subsection{$f_c (l)$ 的一般特征}

总之，利用全部数据，$(f,l)$ 平面中我们穿越线上临界尺寸的粗略估计为 $l_ct_c\approx
0.6$。$l_c\approx \sqrt{\sigma}$，其中 $\sigma$ 是弦张力，我们使用了近似关系 $0.6 \sigma=t_c$。以 QCD 单位计，$l_c$ 约为半个费米，大体符合不断积累的数值数据——后者把从微扰到弦型行为的转变定在相当短、小于一个费米的尺度上。显然，改变穿越线、圈的形状和/或威尔逊圈算符连续统构造的细节会产生略有不同的 $l_c$ 值。
然而，我们发现的临界线 $f_c( l)$ 随标度的小幅变化表明相变将发生在一个受限的尺度范围内。数值上似乎 $f_c(l)$ 当 $l\to 0$ 时渐近到一个正常数。这排除了 $l_c t_c << 1$ 的威尔逊圈定义选择。然而，如果非常细格点上的进一步模拟以某种方式确定了当 $l\to 0$ 时 $f_c(l)\to 0$，那么必须承认这类选择是可能的；
那种情形下，有效弦描述不太可能对太靠近 $l_c$ 的 $l$ 成立。即便可能，我们也不会选择这种受调节连续统威尔逊圈算符的定义。

我们已经看到 $F_c(L,b)$ 对 $f_c (l)$ 的近似在靠近大 $N$ 相变点的一段物理尺度范围内的标度破坏。要把圈的物理尺寸减得更小是不可能的，因为这会要求过大的 $b$ 值，超出大 $N$ 扭转约化技巧在 $7^4$ 格点上可行的范围。类似地，到达大得多的圈也不可能，因为在过小的 $b$ 值下工作会使格点系统经过一级相变进入一个与连续统物理无关的强耦合相。

我们发现 $f_c(l)$ 对小圈变平了。有可能当 $l$ 趋于零时 $f(l)$ 确实降到零，但如果果真发生，可以预期该效应是对数的（即
$f(l) \sim\frac{1}{|\log \Lambda l|}$），因为 $f(l)$ 应当大致表现得像标度 $l$ 处的规范耦合常数，而涂抹旨在消除的周长发散正是由它决定的。当 $l$ 增大时，$f_c(l)$ 开始更快地增长，最终 $f_c(l)$ 将违反其在格点上的界限。要安排太大的圈拥有谱隙是很勉强的，
因为对大圈，禁闭体现在几乎均匀的本征值密度上。

然而请注意，禁闭并非拥有我们的大 $N$ 相变所必需；即使本征值密度对非常大的圈不趋于均匀分布，大圈的谱无隙仍然是可能的。因此，类似的分析或许也适用于不禁闭的模型。
